% ======================================================================
% Informe: Aplicación de SOLID al backend clínico y de catálogo
% Proyecto: tu_vacuna_pai (backend Spring Boot, modulo API)
% Academia: ejercicio de patrones de diseño / Java / Spring
%
% NOTA: se entrega únicamente el fuente LaTeX; no se compila el PDF.
% Compilación sugerida: pdflatex -interaction=nonstopmode
%   solid-backend-informe.tex
% ======================================================================
\documentclass[11pt,a4paper]{article}

\usepackage[utf8]{inputenc}
\usepackage[T1]{fontenc}
\usepackage[spanish]{babel}
\usepackage{geometry}
\usepackage{booktabs}
\usepackage{array}
\usepackage{longtable}
\usepackage{listings}
\usepackage{xcolor}
\usepackage[hidelinks]{hyperref}

\geometry{margin=2.5cm}

\lstset{
  basicstyle=\ttfamily\small,
  breaklines=true,
  columns=fullflexible,
  frame=single,
  rulecolor=\color{gray},
  backgroundcolor=\color{gray!8},
}

\title{\textbf{Aplicación de los principios SOLID al backend}\\
  \large Módulos clínico (\texttt{patients}, \texttt{attentions}),
  catálogo (\texttt{catalog}) y \texttt{shared} \\[4pt]
  \normalsize Proyecto \texttt{tu\_vacuna\_pai} --- Spring Boot 4.1.1 / Java 21}
\author{Equipo de desarrollo --- refactor académico}
\date{12 de septiembre de 2026}

\begin{document}

\maketitle

\tableofcontents
\newpage

% ----------------------------------------------------------------------
\section{Resumen}
% ----------------------------------------------------------------------

Este informe documenta el refactor del backend Spring Boot aplicando los
principios SOLID a los módulos clínico (\texttt{patients},
\texttt{attentions}), al catálogo (\texttt{catalog}) y a utilidades
compartidas (\texttt{shared}). Se trata de un ejercicio académico: el
\textit{frontend} Flutter no se tocó, la arquitectura por capas declarada en
ADR-006 se respeta, y toda verificación se hace sin base de datos
(112 tests unitarios de negocio + 4 reglas ArchUnit = \textbf{116 tests, 0
fallos}).

Los cuatro problemas atacados son:

\begin{enumerate}
  \item \textbf{Inversión de dependencias (DIP)}: el módulo clínico dejó de
        depender de los repositorios JPA del catálogo; hoy consulta un puerto
        propio (\texttt{VaccineCatalogPolicy}).
  \item \textbf{Responsabilidad única (SRP)}: el mapeo entidad\texttt{-}$>$DTO
        se extrajo a mappers (\texttt{AttentionMapper}, \texttt{PatientMapper}).
  \item \textbf{Abierto/cerrado (OCP)}: la ceremonia idempotente y la auditoría
        se encapsularon en \texttt{IdempotencyCoordinator}.
  \item \textbf{DRY y gobierno de permisos}: helpers y chequeos de permiso
        duplicados se centralizaron (\texttt{PermissionGuard}, \texttt{Strings});
        se eliminó código muerto y SQL nativo duplicado.
\end{enumerate}

\section{Contexto y alcance}

El \texttt{AttentionService} (429 líneas) y el \texttt{PatientService} (502
líneas) concentraban varias responsabilidades y dependencias:
\texttt{AttentionService} inyectaba 11 colaboradores, entre ellos cuatro
repositorios del catálogo. \texttt{InstitutionVaccineService} usaba
\texttt{EntityManager} con la misma consulta SQL nativa duplicada para el
upsert ``copy-once'' de habilitación y clonado de catálogo.

Los módulos \texttt{patients} y \texttt{attentions} implementan funcionalidad
completa y probada (\texttt{docs/domain/clinical-module-implementation.md} y
\texttt{vaccine-catalog-implementation.md}), lo que los hace seguros para
refactorizar sin riesgo de regresión funcional (\texttt{mvnw test} como red de
seguridad).

\section{Antes} \label{sec:antes}

A continuación se detallan las violaciones detectadas con su ubicación
(\texttt{archivo:línea}).

\subsection{DIP --- \texttt{AttentionService} depende de la persistencia del catálogo}

\begin{itemize}
  \item \texttt{AttentionService.java:65-75} (antes): inyecta 11 colaboradores,
        entre ellos \texttt{VaccineRepository} (línea 68),
        \texttt{VaccineOptionRepository} (línea 69),
        \texttt{InstitutionVaccineRepository} (línea 70) e
        \texttt{InstitutionVaccineOptionRepository} (línea 71), además de
        \texttt{PatientRepository}.
  \item Validación de catálogo embebida en el servicio:
        \texttt{AttentionService.java:366-404} (antes):
        \texttt{requireEnabledVaccine}, \texttt{requireOption},
        \texttt{resolveInstitutionOption}.
\end{itemize}

\textbf{Consecuencia}: el módulo clínico no podía probarse sin mockear la
persistencia del catálogo, y cualquier cambio en las entidades/repositorios del
catálogo podía romperlo.

\subsection{SRP --- mapeo entidad\texttt{-}$>$DTO dentro de los servicios}

\begin{itemize}
  \item \texttt{AttentionService.java:433-477} (antes):
        \texttt{toResponse}/\texttt{toDoseResponse} (mapeo privado de 45
        líneas).
  \item \texttt{PatientService.java:439-501} (antes):
        \texttt{toResponse} (mapeo privado de 62 líneas, incluye las 5
        subentidades del agregado).
\end{itemize}

\subsection{OCP --- ceremonia idempotente duplicada}

El patrón \texttt{find $\to$ cuerpo $\to$ auditar $\to$ registrar} se repetía en:

\begin{itemize}
  \item \texttt{AttentionService.java:104-139} (antes):
        \texttt{create} ($\approx$35 líneas de ceremonia).
  \item \texttt{AttentionService.java:234-305} (antes):
        \texttt{registerDose} ($\approx$70 líneas mezclando ceremonia y cuerpo).
  \item \texttt{PatientService.java:88-133} (antes):
        \texttt{PatientService.create}.
\end{itemize}

\subsection{DRY y gobierno de permisos}

\begin{itemize}
  \item \texttt{VaccineService.java:29-34} (antes) e
        \texttt{InstitutionVaccineService.java:63-69} (antes):
        método privado \texttt{actor(actorId, permission)} duplicado.
  \item \texttt{AttentionService.java:414-431} y \texttt{PatientService.java:416-437}
        (antes): helpers privados \texttt{blankToNull}, \texttt{safe},
        \texttt{parseOperationId} duplicados.
  \item \texttt{InstitutionVaccineService.java:119-127} y \texttt{205-213}
        (antes): la misma consulta SQL nativa
        \texttt{INSERT \ldots ON CONFLICT \ldots RETURNING id} duplicada en
        \texttt{enable} y \texttt{doClone} usando \texttt{EntityManager}.
  \item \texttt{InstitutionVaccineRepository.java:13-15} (antes):
        \texttt{setEnabled(id, version, enabled)} era código muerto (sin
        llamadores); el commit remoto \texttt{16460fc} ya lo había borrado.
\end{itemize}

\subsection{LSP}

No existen jerarquías de herencia en el alcance; el principio de sustitución de
Liskov no aplica y no introdujo ningún cambio.

\section{Después} \label{sec:despues}

\subsection{DIP: puerto de selección del catálogo}

\textbf{Nueva pieza} \texttt{attentions/service/VaccineCatalogPolicy.java}: un
puerto con un único método

\begin{lstlisting}
DoseSelection resolve(ResolutionRequest request);
\end{lstlisting}

que devuelve un valor inerte (\texttt{DoseSelection}) con el snapshot que la
dosis necesita; \emph{no} expone entidades JPA del catálogo.

\textbf{Nueva pieza} \texttt{catalog/service/CatalogSelectionService.java}:
implementación del puerto que encapsula la validación que antes vivía en el
módulo clínico (vacuna activa y habilitada en la institución, opciones globales
válidas por \texttt{fieldType}, opciones operativas institucionales). El registro
de intereses queda así:

\begin{lstlisting}
AttentionService --usar--> VaccineCatalogPolicy (interfaz, atentions)
CatalogSelectionService (catalog) --implementar--> VaccineCatalogPolicy
\end{lstlisting}

\textbf{Resultado}: \texttt{AttentionService.java:65-75} (después) pasa de 11 a
9 colaboradores y elimina los 4 repositorios de catálogo y las validaciones
listadas en la sección~\ref{sec:antes}. Un \texttt{@ArchTest} nuevo
(\texttt{ArchitectureTest.java}) prohíbe que el módulo \texttt{attentions}
dependa de \texttt{.catalog.repository} o \texttt{.catalog.entity}, y otro exige
que toda implementación del puerto viva en \texttt{catalog}.

\subsection{SRP: mappers de respuesta}

\textbf{Nuevas piezas} \texttt{AttentionMapper.java} (mapeo de
\texttt{AttentionResponse} y \texttt{AppliedDoseResponse}) y
\texttt{PatientMapper.java} (mapeo de \texttt{PatientResponse} con sus cinco
subentidades). Cada servicio inyecta su mapper y conserva solo la
orquestación, la validación de dominio y la persistencia.

\subsection{OCP: coordinador de escrituras idempotentes}

\textbf{Nueva pieza} \texttt{shared/application/IdempotencyCoordinator.java}
encapsula la ceremonia:

\begin{lstlisting}
T execute(operationId, commandType, responseType,
          Supplier<AuditContext> context, Supplier<WriteResult<T>> body) {
  var previous = processed.find(operationId, responseType);
  if (previous.isPresent()) return previous.get();          // replay: no hace nada
  WriteResult<T> result = body.get();                        // dominio del servicio
  audit.record(context.get().actorId(), ..., operationId, result.response());
  processed.record(operationId, commandType, result.aggregateId(), result.response());
  return result.response();
}
\end{lstlisting}

El contexto de auditoría es diferido a propósito: en una repetición idempotente
\textbf{no se ejecuta absolutamente nada}, ni siquiera la resolución del actor.
Se usa \texttt{execute} en \texttt{create} y \texttt{registerDose}, y
\texttt{executeAudited} (sin idempotencia) en \texttt{update},
\texttt{complete}, \texttt{updateContact}, \texttt{updateDemographics} y
\texttt{updateMedicalHistories}. Las cancelaciones conservan su auditoría
específica (\texttt{Cancellation}) dentro de los servicios.

\textbf{Resultado}: la gestión transversal quedó en una sola clase; los
servicios aportan únicamente el cuerpo, cumpliendo el OCP.

\subsection{DRY y gobierno de permisos}

\textbf{Nueva pieza} \texttt{shared/security/PermissionGuard.java}: guardia de
autorización por permiso; sustituye el \texttt{actor()} duplicado en
\texttt{VaccineService} e \texttt{InstitutionVaccineService}.

\textbf{Nueva pieza} \texttt{shared/util/Strings.java}: centraliza
\texttt{blankToNull}, \texttt{trimToNull} y \texttt{safe}.

\textbf{Repositorio}: \texttt{InstitutionVaccineRepository.java} (después)
elimina el \texttt{setEnabled} muerto y declara:

\begin{lstlisting}
int setEnabledById(UUID id, boolean enabled);                    // JPQL con WHERE id
int insertEnabledIfAbsent(UUID institution, UUID vaccine, UUID actor); // INSERT ... ON CONFLICT
\end{lstlisting}

el upsert \texttt{copy-once} regresa 1/0 (filas insertadas), preservando la
semántica del \texttt{RETURNING id} previo sin necesidad de
\texttt{EntityManager}. Con esto \texttt{InstitutionVaccineService.java:34-61}
(después) elimina el \texttt{EntityManager} y el \texttt{IdentityService},
inyectando \texttt{PermissionGuard}.

\section{Comparación antes/después}

\begin{longtable}{@{}p{4.6cm}p{5.2cm}p{5.2cm}@{}}
\toprule
\textbf{Clase} & \textbf{Antes} & \textbf{Después} \\
\midrule
\endhead
\texttt{AttentionService}
 & 429 líneas; 11 colaboradores; 4 repos de \texttt{catalog} +
   \texttt{PatientRepository}; validación de catálogo embebida (líneas
   366-404); mapeo privado (433-477); helpers duplicados (414-431).
 & 336 líneas; 9 colaboradores; consulta vía \texttt{VaccineCatalogPolicy};
   mapeo en \texttt{AttentionMapper}; ceremonia en
   \texttt{IdempotencyCoordinator}; helpers en \texttt{Strings}.
\\
\texttt{PatientService}
 & 502 líneas; 10 colaboradores; mapeo privado de 62 líneas (439-501);
   helpers \texttt{blankToNull}/\texttt{safe}/\texttt{parseOperationId}
   (416-437).
 & 433 líneas; mapeo en \texttt{PatientMapper}; ceremonia en
   \texttt{IdempotencyCoordinator}; helpers en \texttt{Strings}.
\\
\texttt{VaccineService}
 & \texttt{IdentityService} + \texttt{actor()} privado (29-34); lógica
   \texttt{isDefault} y \texttt{lockActive} en \texttt{updateOption}/
   \texttt{updateTemplate}.
 & \texttt{PermissionGuard}; misma lógica de dominio.
\\
\texttt{InstitutionVaccineService}
 & \texttt{IdentityService} + \texttt{actor()} (63-69); \texttt{EntityManager}
   con SQL nativo duplicado en \texttt{enable} (119-127) y \texttt{doClone}
   (205-213).
 & \texttt{PermissionGuard}; upsert declarado en el repositorio;
   sin \texttt{EntityManager}.
\\
\texttt{InstitutionVaccineRepository}
 & \texttt{setEnabled(id, version, enabled)} muerto (13-15).
 & \texttt{setEnabledById} + \texttt{insertEnabledIfAbsent} (nativo
   copy-once).
\\
\bottomrule
\end{longtable}

\section{Arquitectura resultante}

\begin{itemize}
  \item \texttt{attentions} $\to$ \texttt{VaccineCatalogPolicy} (puerto).
  \item \texttt{catalog} $\to$ \texttt{CatalogSelectionService} (implementa el
        puerto; es la única dependencia entre módulos, fuera de \texttt{shared}).
  \item \texttt{shared.application} $\to$ \texttt{IdempotencyCoordinator}.
  \item \texttt{shared.security} $\to$ \texttt{PermissionGuard}.
  \item \texttt{shared.util} $\to$ \texttt{Strings}.
  \item Reglas ArchUnit nuevas en
        \texttt{shared/architecture/ArchitectureTest.java}: frontera
        \texttt{attentions} {$\nrightarrow$} \texttt{catalog.repository/entity}
        y radicación del puerto en \texttt{catalog}.
\end{itemize}

\section{Verificación}

\begin{lstlisting}
mvnw.cmd compile && mvnw.cmd test
\end{lstlisting}

\begin{itemize}
  \item Sin base de datos: todos los tests son JUnit 5 + Mockito puros.
  \item Línea base: 114 tests, 0 fallos.
  \item Tras el refactor: \textbf{116 tests, 0 fallos} (se añaden 2 reglas
        ArchUnit). \texttt{BUILD SUCCESS}.
  \item Tests actualizados: \texttt{AttentionServiceTest},
        \texttt{PatientServiceTest}, \texttt{InstitutionVaccineServiceCloneTest}
        (ya no mockean \texttt{EntityManager} ni repositorios de catálogo;
        mockean los nuevos colaboradores).
  \item El backend también compila y arranca con \texttt{spring-boot:run}
        (health \texttt{UP}) usando el JDK real
        \texttt{C:\textbackslash Program Files\textbackslash Java\textbackslash jdk-26.0.2.1}.
\end{itemize}

\section{Recomendaciones futuras}

\begin{enumerate}
  \item \textbf{Integrar \texttt{UserProvisioningService}} con
        \texttt{PermissionGuard} y \texttt{Strings} (quedó fuera del alcance por
        riesgo; ya tiene los recursos listos).
  \item \textbf{Cancelaciones en el coordinador}: añadir un overload de
        \texttt{executeAudited} con \texttt{payloadFn} para llevar
        \texttt{cancel}/\texttt{cancelDose} a la misma ceremonia sin perder su
        payload de auditoría compuesto.
  \item \textbf{Evitar la doble resolución del actor} en cuerpos que ya lo
        resuelven: evaluar que \texttt{AuditContext} transporte el
        \texttt{AuthorizedUser} en lugar de solo el id.
  \item \textbf{Dependencia residual}: \texttt{AttentionService} aún importa
        \texttt{catalog.exception.OptimisticCatalogException}; migrar esa
        excepción genérica de versión optimista a \texttt{shared.exceptions}.
  \item \textbf{Tests}: cubrir \texttt{CatalogSelectionService} con casos de
        vacuna inactiva, no habilitada y opciones inválidas; cubrir
        \texttt{PermissionGuard}.
  \item \textbf{Mappers}: si el volumen de mapeo crece, evaluar MapStruct en
        lugar de mappers manuales (sin cambiar la estructura expuesta).
  \item \textbf{Sync}: cuando se implemente el \texttt{SyncEngine}, reutilizar
        \texttt{IdempotencyCoordinator} para que el camino \texttt{/sync/push}
        comparta la misma garantía D3.
\end{enumerate}

\end{document}